% ======================================================================
%  AT-MONO102 — Chapter 3: Actualization
%  Canonical Monograph (v2.0) · The Actualization Theory:
%  A Reconstruction of Physics from Difference, Actualization and Spectrum
%
%  Status: COMPLETE — PASS-READY (MONO007 referee-ready architecture)
%  Inputs: Chapter01_Difference.tex, Chapter02_Eta.tex,
%          ATQG_CanonicalMonograph.md, ATQG_Mono007RefereeReadiness.md
%  Method: assembly only — canonical results only, no new physics, no speculation
%  Citations: QG268 (count conservation), QG272 (sector emergence),
%             QG282 (closure principle), QG292 (minimal foundation),
%             QG293 (hierarchy necessity), QG294 (minimal theory),
%             MONO006 (actualization derived), MONO007 (referee readiness)
%
%  Usage: % ======================================================================
%  AT-MONO102 — Chapter 3: Actualization
%  Canonical Monograph (v2.0) · The Actualization Theory:
%  A Reconstruction of Physics from Difference, Actualization and Spectrum
%
%  Status: COMPLETE — PASS-READY (MONO007 referee-ready architecture)
%  Inputs: Chapter01_Difference.tex, Chapter02_Eta.tex,
%          ATQG_CanonicalMonograph.md, ATQG_Mono007RefereeReadiness.md
%  Method: assembly only — canonical results only, no new physics, no speculation
%  Citations: QG268 (count conservation), QG272 (sector emergence),
%             QG282 (closure principle), QG292 (minimal foundation),
%             QG293 (hierarchy necessity), QG294 (minimal theory),
%             MONO006 (actualization derived), MONO007 (referee readiness)
%
%  Usage: % ======================================================================
%  AT-MONO102 — Chapter 3: Actualization
%  Canonical Monograph (v2.0) · The Actualization Theory:
%  A Reconstruction of Physics from Difference, Actualization and Spectrum
%
%  Status: COMPLETE — PASS-READY (MONO007 referee-ready architecture)
%  Inputs: Chapter01_Difference.tex, Chapter02_Eta.tex,
%          ATQG_CanonicalMonograph.md, ATQG_Mono007RefereeReadiness.md
%  Method: assembly only — canonical results only, no new physics, no speculation
%  Citations: QG268 (count conservation), QG272 (sector emergence),
%             QG282 (closure principle), QG292 (minimal foundation),
%             QG293 (hierarchy necessity), QG294 (minimal theory),
%             MONO006 (actualization derived), MONO007 (referee readiness)
%
%  Usage: % ======================================================================
%  AT-MONO102 — Chapter 3: Actualization
%  Canonical Monograph (v2.0) · The Actualization Theory:
%  A Reconstruction of Physics from Difference, Actualization and Spectrum
%
%  Status: COMPLETE — PASS-READY (MONO007 referee-ready architecture)
%  Inputs: Chapter01_Difference.tex, Chapter02_Eta.tex,
%          ATQG_CanonicalMonograph.md, ATQG_Mono007RefereeReadiness.md
%  Method: assembly only — canonical results only, no new physics, no speculation
%  Citations: QG268 (count conservation), QG272 (sector emergence),
%             QG282 (closure principle), QG292 (minimal foundation),
%             QG293 (hierarchy necessity), QG294 (minimal theory),
%             MONO006 (actualization derived), MONO007 (referee readiness)
%
%  Usage: \input{Chapter03_Actualization.tex} after the preamble that defines
%         the theorem environments below (or include via the master file).
% ======================================================================

% ---- Theorem environments (define in the master preamble if not present) ----
% \newtheorem{definition}{Definition}[chapter]
% \newtheorem{lemma}[definition]{Lemma}
% \newtheorem{theorem}[definition]{Theorem}
% \newtheorem{corollary}[definition]{Corollary}
% \newtheorem{remark}[definition]{Remark}

\chapter{Actualization}
\label{c03:chap:actualization}

\section{Introduction}
\label{c03:sec:actualization-introduction}

In Chapter~\ref{c01:chap:difference} we established that Difference is the fundamental boundary
of the theory --- the primitive whose count-conservation law is its definitional identity ---
and in Chapter~\ref{c02:chap:eta} that the tensor reference $\eta$ is the second primitive,
against which the tensor face of Difference is read. The canonical hierarchy
\eqref{c03:eq:core} then proceeds from the primitive pair to the inevitable spectrum through a
process:

\begin{equation}
\label{c03:eq:core}
\text{Difference}\;\longrightarrow\;\text{Actualization}\;\longrightarrow\;\text{Inevitable Spectrum}\;\longrightarrow\;\text{Physics}.
\end{equation}

This chapter develops the second member of that hierarchy: \emph{Actualization}. We
establish, with formal precision, that Actualization is the \emph{process face of
Difference} --- the count-producing dynamics whose definitional identity is count
conservation \cite{c03:QG268}; that its dynamics converge to the stable fixed point $N=96$, the
closure of the theory \cite{c03:QG282}; that the resonance structure of the spectrum is
actualization's readout \cite{c03:QG272}; and --- decisively, consistent with the MONO006
resolution \cite{c03:MONO006} --- that Actualization is a \emph{derived necessity}, not a third
primitive.

\section{Actualization as the Process Face of Difference}
\label{c03:sec:process-face}

\begin{theorem}[Actualization is the process face of Difference]
\label{c03:thm:process-face}
\emph{Actualization} --- the count-producing dynamics by which a Q-event, the unit of
Difference \cite{c03:QG268}, is applied, generating the network whose convergence yields the
spectrum --- is not an independent entity: it is the process face of the one primitive
Difference. Consequently the primitives of the theory are exactly
$\{\text{Difference},\eta\}$: Actualization is their derived process face, not a primitive.
\end{theorem}

\begin{proof}
Count conservation is the definitional identity of the primitive \cite{c03:QG268}: Difference is
the counting difference from a uniform background, and a Q-event is its unit. Actualization
is precisely the process that applies those units, and its identity is inherited from the
primitive: there is no Actualization without the unit, and no unit without Difference. The
minimal primitive set was established as the pair $\{\text{Difference},\eta\}$
(Theorem~\ref{c02:thm:minimal-foundation}, Chapter~\ref{c02:chap:eta}); the MONO006 resolution
\cite{c03:MONO006} confirms this classification, and the MONO007 architecture \cite{c03:MONO007}
presents Actualization as a derived layer. Hence there is no primitive-status ambiguity:
Actualization is the dynamical face of Difference, not a separate primitive.
\end{proof}

\section{Count Conservation}
\label{c03:sec:count-conservation}

\begin{theorem}[Count conservation is the definitional identity]
\label{c03:thm:count-conservation}
\emph{Count conservation} --- the total count of units of Difference is conserved by the
actualization process: Q-events are neither created nor destroyed, only redistributed ---
is not a separate postulate but the definitional identity of the primitive Difference and
therefore of its process face Actualization. One application of Actualization produces one
unit of count, so conservation follows automatically from the primitive, not as an imported
assumption of the theory.
\end{theorem}

\begin{proof}
The count-conservation origin \cite{c03:QG268} establishes that conservation is the defining
property of the primitive itself --- it is what \emph{is} the count --- and identifies the
Q-event as the indivisible step of the count. Since Actualization is the process that
applies the count (Theorem~\ref{c03:thm:process-face}), the conservation law is inherited:
the process conserves exactly what the primitive defines. No separate conservation postulate
is required.
\end{proof}

\section{The \texorpdfstring{$N=96$}{N=96} Attractor}
\label{c03:sec:n96-attractor}

The \emph{closure} of the theory is the stable fixed point of the actualization process: the
configuration at which the count-producing dynamics has converged, with zero residual
change. The boundary of the theory is this fixed point --- the topology saturates at zero
residual link growth, and every initial pattern converges to the same final geometry. This
is the fixed point that Chapter~\ref{c01:chap:difference} invoked in characterizing (not
deriving) the fundamental boundary (Corollary~\ref{c01:cor:closure-characterizes}); here we
see its origin: the fixed point is the convergence of Difference's own count-producing
process.

\begin{theorem}[The attractor is \texorpdfstring{$N=96$}{N=96}]
\label{c03:thm:n96}
The converged attractor of the actualization process is the $N=96$ network: the size $N=96$
is the stable fixed point of the dynamics, not a freely chosen input.
\end{theorem}

\begin{proof}
The boundary-origin audit (the Closure Principle \cite{c03:QG282}) establishes that the
observable sector is the converged attractor of the actualization dynamics. The dynamics has
a self-reinforcing feedback (activity to links to activity) bounded by network capacity, so
positive feedback saturates at a stable fixed point: zero residual link growth and a
content-independent final geometry. The convergence to the $N=96$ attractor follows from the
closure principle \cite{c03:QG282} and the reduction program \cite{c03:QG293,c03:QG294};
$N=96$ is thereby an output of the dynamics --- the inevitable fixed point --- not a
primitive or a free parameter.
\end{proof}

\begin{corollary}[The spectrum is the attractor's eigenspectrum]
\label{c03:cor:spectrum-eigenspectrum}
The D96 spectrum is the Laplacian eigenspectrum of the converged $N=96$ network: the
inevitable output of the actualization attractor, not an additional primitive.
\end{corollary}

\begin{proof}
The converged network defines a graph Laplacian; its eigenspectrum is the D96 spectrum, with
the 95 positive modes and the single zero mode (the background). Since the network is the
converged fixed point of the dynamics (Theorem~\ref{c03:thm:n96}), the spectrum is a
deterministic function of the attractor --- an output, not an input. This is the inevitable-
spectrum result of the canonical program \cite{c03:QG293,c03:QG294}.
\end{proof}

\section{Actualization and Resonance}
\label{c03:sec:actualization-resonance}

\begin{theorem}[Resonance is actualization's readout]
\label{c03:thm:resonance-readout}
\emph{Resonance} --- the spectral readout of the actualization process, the operator layer
projecting the converged spectrum onto the observable structure --- is actualization's
readout: the sector structure of observable physics is a projection of the one
actualization-driven spectrum. The resonance condition of the theory is the conjunction of
count conservation (the definitional identity) and the closure boundary (the fixed point):
Resonance equals Conservation plus Boundary.
\end{theorem}

\begin{proof}
The sector-emergence audit \cite{c03:QG272} establishes that distinct sectors --- masses,
couplings, mixings, gravity, cosmology --- are projection classes of a single operator layer
over the one spectrum, with no sector-exclusive operator. Since the spectrum is the output
of the actualization attractor (Corollary~\ref{c03:cor:spectrum-eigenspectrum}), the resonance
structure read from it is the readout of the actualization process: one dynamics, one
spectrum, one operator basis, projected onto distinct physical questions. Conservation is
the definitional identity of the count (Theorem~\ref{c03:thm:count-conservation}); the
boundary is the closure fixed point $N=96$ (Theorem~\ref{c03:thm:n96}); the spectrum is
determined by the process that conserves its count and closes at its boundary. The resonance
condition is therefore the conjunction of the two --- Conservation plus Boundary ---
actualization's spectral face.
\end{proof}

\section{Actualization as Derived Necessity}
\label{c03:sec:derived-necessity}

\begin{theorem}[Actualization is derived, not primitive]
\label{c03:thm:derived-not-primitive}
Actualization is INDISPENSABLE as a derivation layer --- without it there is no count, no
network, no spectrum, and no physics --- but it is a DERIVED necessity: indispensable as a
derivation step, yet fully dependent on the primitive Difference and therefore not an
underivable input.
\end{theorem}

\begin{proof}
The hierarchy-necessity audit \cite{c03:QG293} classifies Actualization as INDISPENSABLE: it is
the count-producing process, so without it there is no count, no network, no spectrum, and
no physics. But necessity as a step is not primitiveness as an input: by
Theorem~\ref{c03:thm:process-face}, Actualization is the process face of Difference; the
foundation stress test \cite{c03:QG292} establishes that removing Difference collapses
Actualization (it needs the definitional unit) while removing $\eta$ leaves it intact
(counting needs no metric reference). The minimal-theory audit \cite{c03:QG294} confirms that
the minimal hierarchy Difference → Actualization → Spectrum → Physics is complete, with
Actualization an essential derived step; the MONO006 resolution \cite{c03:MONO006} concludes:
Actualization is derived from Difference alone, indispensable as a step, not a third
primitive.
\end{proof}

The canonical architecture is therefore unchanged --- primitives $\{\text{Difference},\eta\}$,
Actualization derived from Difference --- consistent with Chapters~\ref{c01:chap:difference}
and \ref{c02:chap:eta}; consistent with MONO007 \cite{c03:MONO007}, all statements that could
be read as ascribing primitive status to Actualization are worded to make the derivation
explicit, and no ambiguity remains.

\section{Consequences}
\label{c03:sec:actualization-consequences}

\begin{theorem}[The spectrum is the inevitable output]
\label{c03:thm:spectrum-inevitable}
The D96 spectrum is the inevitable output of the actualization attractor: a deterministic
function of the converged network, carrying no independent choice.
\end{theorem}

\begin{proof}
By Corollary~\ref{c03:cor:spectrum-eigenspectrum}, the spectrum is the Laplacian eigenspectrum of
the converged $N=96$ attractor. The attractor is unique and content-independent
(Theorem~\ref{c03:thm:n96}); the eigenspectrum is a deterministic function of it. Hence the
spectrum carries no choice and is inevitable --- the output of the actualization process,
not an input to it. The reduction program over the fuller hierarchy thereby closes at the
minimal hierarchy Difference → Actualization → Spectrum → Physics: the hierarchy-necessity
audit \cite{c03:QG293} classifies the intermediate layers --- Closure COMPRESSIBLE into
Actualization (the fixed point of the process), Conservation COMPRESSIBLE into
Difference/network (the definitional identity plus the graph handshake), Resonance
COMPRESSIBLE into Spectrum (the operator readout) --- and the minimal-theory audit
\cite{c03:QG294} verifies that every compressed layer is derivable from the minimal hierarchy,
with none actually required as an input.
\end{proof}

\begin{corollary}[Physics is the readout of the actualized spectrum]
\label{c03:cor:physics-readout}
All observable physics is the readout of the actualized spectrum: the resonance structure
projected by the operator layer, from which the sectors of physics emerge.
\end{corollary}

\section{Summary}
\label{c03:sec:actualization-summary}

This chapter has established the second member of the canonical architecture. Actualization
is the process face of Difference (Theorem~\ref{c03:thm:process-face}), whose count
conservation is the primitive's definitional identity (Theorem~\ref{c03:thm:count-conservation})
and whose dynamics converge to the stable fixed point $N=96$
(Theorem~\ref{c03:thm:n96}, Corollary~\ref{c03:cor:spectrum-eigenspectrum}); the resonance
structure read from the resulting spectrum satisfies Resonance equals Conservation plus
Boundary (Theorem~\ref{c03:thm:resonance-readout}). Actualization is indispensable as a layer
yet derived, not primitive (Theorem~\ref{c03:thm:derived-not-primitive}), and the reduction
closes at the minimal hierarchy, with physics the readout of the actualized spectrum
(Theorem~\ref{c03:thm:spectrum-inevitable}, Corollary~\ref{c03:cor:physics-readout}).

The following chapter derives, from this process, the spectrum itself.

\begin{thebibliography}{99}
\bibitem{c03:QG268} AT-QG Phase 268, \emph{Count Conservation Origin}, Docs/Research.
\bibitem{c03:QG272} AT-QG Phase 272, \emph{Sector Emergence Audit}, Docs/Research.
\bibitem{c03:QG282} AT-QG Phase 282, \emph{Boundary Origin Audit} (Closure Principle), Docs/Research.
\bibitem{c03:QG292} AT-QG Phase 292, \emph{Foundation Stress Test}, Docs/Research.
\bibitem{c03:QG293} AT-QG Phase 293, \emph{Hierarchy Necessity Audit}, Docs/Research.
\bibitem{c03:QG294} AT-QG Phase 294, \emph{Minimal Theory Audit}, Docs/Research.
\bibitem{c03:MONO006} AT-MONO006, \emph{A01 Actualization Origin Audit}, Docs/Zenodo/Monograph.
\bibitem{c03:MONO007} AT-MONO007, \emph{Referee-Readiness Resolution}, Docs/Zenodo/Monograph.
\end{thebibliography} after the preamble that defines
%         the theorem environments below (or include via the master file).
% ======================================================================

% ---- Theorem environments (define in the master preamble if not present) ----
% \newtheorem{definition}{Definition}[chapter]
% \newtheorem{lemma}[definition]{Lemma}
% \newtheorem{theorem}[definition]{Theorem}
% \newtheorem{corollary}[definition]{Corollary}
% \newtheorem{remark}[definition]{Remark}

\chapter{Actualization}
\label{c03:chap:actualization}

\section{Introduction}
\label{c03:sec:actualization-introduction}

In Chapter~\ref{c01:chap:difference} we established that Difference is the fundamental boundary
of the theory --- the primitive whose count-conservation law is its definitional identity ---
and in Chapter~\ref{c02:chap:eta} that the tensor reference $\eta$ is the second primitive,
against which the tensor face of Difference is read. The canonical hierarchy
\eqref{c03:eq:core} then proceeds from the primitive pair to the inevitable spectrum through a
process:

\begin{equation}
\label{c03:eq:core}
\text{Difference}\;\longrightarrow\;\text{Actualization}\;\longrightarrow\;\text{Inevitable Spectrum}\;\longrightarrow\;\text{Physics}.
\end{equation}

This chapter develops the second member of that hierarchy: \emph{Actualization}. We
establish, with formal precision, that Actualization is the \emph{process face of
Difference} --- the count-producing dynamics whose definitional identity is count
conservation \cite{c03:QG268}; that its dynamics converge to the stable fixed point $N=96$, the
closure of the theory \cite{c03:QG282}; that the resonance structure of the spectrum is
actualization's readout \cite{c03:QG272}; and --- decisively, consistent with the MONO006
resolution \cite{c03:MONO006} --- that Actualization is a \emph{derived necessity}, not a third
primitive.

\section{Actualization as the Process Face of Difference}
\label{c03:sec:process-face}

\begin{theorem}[Actualization is the process face of Difference]
\label{c03:thm:process-face}
\emph{Actualization} --- the count-producing dynamics by which a Q-event, the unit of
Difference \cite{c03:QG268}, is applied, generating the network whose convergence yields the
spectrum --- is not an independent entity: it is the process face of the one primitive
Difference. Consequently the primitives of the theory are exactly
$\{\text{Difference},\eta\}$: Actualization is their derived process face, not a primitive.
\end{theorem}

\begin{proof}
Count conservation is the definitional identity of the primitive \cite{c03:QG268}: Difference is
the counting difference from a uniform background, and a Q-event is its unit. Actualization
is precisely the process that applies those units, and its identity is inherited from the
primitive: there is no Actualization without the unit, and no unit without Difference. The
minimal primitive set was established as the pair $\{\text{Difference},\eta\}$
(Theorem~\ref{c02:thm:minimal-foundation}, Chapter~\ref{c02:chap:eta}); the MONO006 resolution
\cite{c03:MONO006} confirms this classification, and the MONO007 architecture \cite{c03:MONO007}
presents Actualization as a derived layer. Hence there is no primitive-status ambiguity:
Actualization is the dynamical face of Difference, not a separate primitive.
\end{proof}

\section{Count Conservation}
\label{c03:sec:count-conservation}

\begin{theorem}[Count conservation is the definitional identity]
\label{c03:thm:count-conservation}
\emph{Count conservation} --- the total count of units of Difference is conserved by the
actualization process: Q-events are neither created nor destroyed, only redistributed ---
is not a separate postulate but the definitional identity of the primitive Difference and
therefore of its process face Actualization. One application of Actualization produces one
unit of count, so conservation follows automatically from the primitive, not as an imported
assumption of the theory.
\end{theorem}

\begin{proof}
The count-conservation origin \cite{c03:QG268} establishes that conservation is the defining
property of the primitive itself --- it is what \emph{is} the count --- and identifies the
Q-event as the indivisible step of the count. Since Actualization is the process that
applies the count (Theorem~\ref{c03:thm:process-face}), the conservation law is inherited:
the process conserves exactly what the primitive defines. No separate conservation postulate
is required.
\end{proof}

\section{The \texorpdfstring{$N=96$}{N=96} Attractor}
\label{c03:sec:n96-attractor}

The \emph{closure} of the theory is the stable fixed point of the actualization process: the
configuration at which the count-producing dynamics has converged, with zero residual
change. The boundary of the theory is this fixed point --- the topology saturates at zero
residual link growth, and every initial pattern converges to the same final geometry. This
is the fixed point that Chapter~\ref{c01:chap:difference} invoked in characterizing (not
deriving) the fundamental boundary (Corollary~\ref{c01:cor:closure-characterizes}); here we
see its origin: the fixed point is the convergence of Difference's own count-producing
process.

\begin{theorem}[The attractor is \texorpdfstring{$N=96$}{N=96}]
\label{c03:thm:n96}
The converged attractor of the actualization process is the $N=96$ network: the size $N=96$
is the stable fixed point of the dynamics, not a freely chosen input.
\end{theorem}

\begin{proof}
The boundary-origin audit (the Closure Principle \cite{c03:QG282}) establishes that the
observable sector is the converged attractor of the actualization dynamics. The dynamics has
a self-reinforcing feedback (activity to links to activity) bounded by network capacity, so
positive feedback saturates at a stable fixed point: zero residual link growth and a
content-independent final geometry. The convergence to the $N=96$ attractor follows from the
closure principle \cite{c03:QG282} and the reduction program \cite{c03:QG293,c03:QG294};
$N=96$ is thereby an output of the dynamics --- the inevitable fixed point --- not a
primitive or a free parameter.
\end{proof}

\begin{corollary}[The spectrum is the attractor's eigenspectrum]
\label{c03:cor:spectrum-eigenspectrum}
The D96 spectrum is the Laplacian eigenspectrum of the converged $N=96$ network: the
inevitable output of the actualization attractor, not an additional primitive.
\end{corollary}

\begin{proof}
The converged network defines a graph Laplacian; its eigenspectrum is the D96 spectrum, with
the 95 positive modes and the single zero mode (the background). Since the network is the
converged fixed point of the dynamics (Theorem~\ref{c03:thm:n96}), the spectrum is a
deterministic function of the attractor --- an output, not an input. This is the inevitable-
spectrum result of the canonical program \cite{c03:QG293,c03:QG294}.
\end{proof}

\section{Actualization and Resonance}
\label{c03:sec:actualization-resonance}

\begin{theorem}[Resonance is actualization's readout]
\label{c03:thm:resonance-readout}
\emph{Resonance} --- the spectral readout of the actualization process, the operator layer
projecting the converged spectrum onto the observable structure --- is actualization's
readout: the sector structure of observable physics is a projection of the one
actualization-driven spectrum. The resonance condition of the theory is the conjunction of
count conservation (the definitional identity) and the closure boundary (the fixed point):
Resonance equals Conservation plus Boundary.
\end{theorem}

\begin{proof}
The sector-emergence audit \cite{c03:QG272} establishes that distinct sectors --- masses,
couplings, mixings, gravity, cosmology --- are projection classes of a single operator layer
over the one spectrum, with no sector-exclusive operator. Since the spectrum is the output
of the actualization attractor (Corollary~\ref{c03:cor:spectrum-eigenspectrum}), the resonance
structure read from it is the readout of the actualization process: one dynamics, one
spectrum, one operator basis, projected onto distinct physical questions. Conservation is
the definitional identity of the count (Theorem~\ref{c03:thm:count-conservation}); the
boundary is the closure fixed point $N=96$ (Theorem~\ref{c03:thm:n96}); the spectrum is
determined by the process that conserves its count and closes at its boundary. The resonance
condition is therefore the conjunction of the two --- Conservation plus Boundary ---
actualization's spectral face.
\end{proof}

\section{Actualization as Derived Necessity}
\label{c03:sec:derived-necessity}

\begin{theorem}[Actualization is derived, not primitive]
\label{c03:thm:derived-not-primitive}
Actualization is INDISPENSABLE as a derivation layer --- without it there is no count, no
network, no spectrum, and no physics --- but it is a DERIVED necessity: indispensable as a
derivation step, yet fully dependent on the primitive Difference and therefore not an
underivable input.
\end{theorem}

\begin{proof}
The hierarchy-necessity audit \cite{c03:QG293} classifies Actualization as INDISPENSABLE: it is
the count-producing process, so without it there is no count, no network, no spectrum, and
no physics. But necessity as a step is not primitiveness as an input: by
Theorem~\ref{c03:thm:process-face}, Actualization is the process face of Difference; the
foundation stress test \cite{c03:QG292} establishes that removing Difference collapses
Actualization (it needs the definitional unit) while removing $\eta$ leaves it intact
(counting needs no metric reference). The minimal-theory audit \cite{c03:QG294} confirms that
the minimal hierarchy Difference → Actualization → Spectrum → Physics is complete, with
Actualization an essential derived step; the MONO006 resolution \cite{c03:MONO006} concludes:
Actualization is derived from Difference alone, indispensable as a step, not a third
primitive.
\end{proof}

The canonical architecture is therefore unchanged --- primitives $\{\text{Difference},\eta\}$,
Actualization derived from Difference --- consistent with Chapters~\ref{c01:chap:difference}
and \ref{c02:chap:eta}; consistent with MONO007 \cite{c03:MONO007}, all statements that could
be read as ascribing primitive status to Actualization are worded to make the derivation
explicit, and no ambiguity remains.

\section{Consequences}
\label{c03:sec:actualization-consequences}

\begin{theorem}[The spectrum is the inevitable output]
\label{c03:thm:spectrum-inevitable}
The D96 spectrum is the inevitable output of the actualization attractor: a deterministic
function of the converged network, carrying no independent choice.
\end{theorem}

\begin{proof}
By Corollary~\ref{c03:cor:spectrum-eigenspectrum}, the spectrum is the Laplacian eigenspectrum of
the converged $N=96$ attractor. The attractor is unique and content-independent
(Theorem~\ref{c03:thm:n96}); the eigenspectrum is a deterministic function of it. Hence the
spectrum carries no choice and is inevitable --- the output of the actualization process,
not an input to it. The reduction program over the fuller hierarchy thereby closes at the
minimal hierarchy Difference → Actualization → Spectrum → Physics: the hierarchy-necessity
audit \cite{c03:QG293} classifies the intermediate layers --- Closure COMPRESSIBLE into
Actualization (the fixed point of the process), Conservation COMPRESSIBLE into
Difference/network (the definitional identity plus the graph handshake), Resonance
COMPRESSIBLE into Spectrum (the operator readout) --- and the minimal-theory audit
\cite{c03:QG294} verifies that every compressed layer is derivable from the minimal hierarchy,
with none actually required as an input.
\end{proof}

\begin{corollary}[Physics is the readout of the actualized spectrum]
\label{c03:cor:physics-readout}
All observable physics is the readout of the actualized spectrum: the resonance structure
projected by the operator layer, from which the sectors of physics emerge.
\end{corollary}

\section{Summary}
\label{c03:sec:actualization-summary}

This chapter has established the second member of the canonical architecture. Actualization
is the process face of Difference (Theorem~\ref{c03:thm:process-face}), whose count
conservation is the primitive's definitional identity (Theorem~\ref{c03:thm:count-conservation})
and whose dynamics converge to the stable fixed point $N=96$
(Theorem~\ref{c03:thm:n96}, Corollary~\ref{c03:cor:spectrum-eigenspectrum}); the resonance
structure read from the resulting spectrum satisfies Resonance equals Conservation plus
Boundary (Theorem~\ref{c03:thm:resonance-readout}). Actualization is indispensable as a layer
yet derived, not primitive (Theorem~\ref{c03:thm:derived-not-primitive}), and the reduction
closes at the minimal hierarchy, with physics the readout of the actualized spectrum
(Theorem~\ref{c03:thm:spectrum-inevitable}, Corollary~\ref{c03:cor:physics-readout}).

The following chapter derives, from this process, the spectrum itself.

\begin{thebibliography}{99}
\bibitem{c03:QG268} AT-QG Phase 268, \emph{Count Conservation Origin}, Docs/Research.
\bibitem{c03:QG272} AT-QG Phase 272, \emph{Sector Emergence Audit}, Docs/Research.
\bibitem{c03:QG282} AT-QG Phase 282, \emph{Boundary Origin Audit} (Closure Principle), Docs/Research.
\bibitem{c03:QG292} AT-QG Phase 292, \emph{Foundation Stress Test}, Docs/Research.
\bibitem{c03:QG293} AT-QG Phase 293, \emph{Hierarchy Necessity Audit}, Docs/Research.
\bibitem{c03:QG294} AT-QG Phase 294, \emph{Minimal Theory Audit}, Docs/Research.
\bibitem{c03:MONO006} AT-MONO006, \emph{A01 Actualization Origin Audit}, Docs/Zenodo/Monograph.
\bibitem{c03:MONO007} AT-MONO007, \emph{Referee-Readiness Resolution}, Docs/Zenodo/Monograph.
\end{thebibliography} after the preamble that defines
%         the theorem environments below (or include via the master file).
% ======================================================================

% ---- Theorem environments (define in the master preamble if not present) ----
% \newtheorem{definition}{Definition}[chapter]
% \newtheorem{lemma}[definition]{Lemma}
% \newtheorem{theorem}[definition]{Theorem}
% \newtheorem{corollary}[definition]{Corollary}
% \newtheorem{remark}[definition]{Remark}

\chapter{Actualization}
\label{c03:chap:actualization}

\section{Introduction}
\label{c03:sec:actualization-introduction}

In Chapter~\ref{c01:chap:difference} we established that Difference is the fundamental boundary
of the theory --- the primitive whose count-conservation law is its definitional identity ---
and in Chapter~\ref{c02:chap:eta} that the tensor reference $\eta$ is the second primitive,
against which the tensor face of Difference is read. The canonical hierarchy
\eqref{c03:eq:core} then proceeds from the primitive pair to the inevitable spectrum through a
process:

\begin{equation}
\label{c03:eq:core}
\text{Difference}\;\longrightarrow\;\text{Actualization}\;\longrightarrow\;\text{Inevitable Spectrum}\;\longrightarrow\;\text{Physics}.
\end{equation}

This chapter develops the second member of that hierarchy: \emph{Actualization}. We
establish, with formal precision, that Actualization is the \emph{process face of
Difference} --- the count-producing dynamics whose definitional identity is count
conservation \cite{c03:QG268}; that its dynamics converge to the stable fixed point $N=96$, the
closure of the theory \cite{c03:QG282}; that the resonance structure of the spectrum is
actualization's readout \cite{c03:QG272}; and --- decisively, consistent with the MONO006
resolution \cite{c03:MONO006} --- that Actualization is a \emph{derived necessity}, not a third
primitive.

\section{Actualization as the Process Face of Difference}
\label{c03:sec:process-face}

\begin{theorem}[Actualization is the process face of Difference]
\label{c03:thm:process-face}
\emph{Actualization} --- the count-producing dynamics by which a Q-event, the unit of
Difference \cite{c03:QG268}, is applied, generating the network whose convergence yields the
spectrum --- is not an independent entity: it is the process face of the one primitive
Difference. Consequently the primitives of the theory are exactly
$\{\text{Difference},\eta\}$: Actualization is their derived process face, not a primitive.
\end{theorem}

\begin{proof}
Count conservation is the definitional identity of the primitive \cite{c03:QG268}: Difference is
the counting difference from a uniform background, and a Q-event is its unit. Actualization
is precisely the process that applies those units, and its identity is inherited from the
primitive: there is no Actualization without the unit, and no unit without Difference. The
minimal primitive set was established as the pair $\{\text{Difference},\eta\}$
(Theorem~\ref{c02:thm:minimal-foundation}, Chapter~\ref{c02:chap:eta}); the MONO006 resolution
\cite{c03:MONO006} confirms this classification, and the MONO007 architecture \cite{c03:MONO007}
presents Actualization as a derived layer. Hence there is no primitive-status ambiguity:
Actualization is the dynamical face of Difference, not a separate primitive.
\end{proof}

\section{Count Conservation}
\label{c03:sec:count-conservation}

\begin{theorem}[Count conservation is the definitional identity]
\label{c03:thm:count-conservation}
\emph{Count conservation} --- the total count of units of Difference is conserved by the
actualization process: Q-events are neither created nor destroyed, only redistributed ---
is not a separate postulate but the definitional identity of the primitive Difference and
therefore of its process face Actualization. One application of Actualization produces one
unit of count, so conservation follows automatically from the primitive, not as an imported
assumption of the theory.
\end{theorem}

\begin{proof}
The count-conservation origin \cite{c03:QG268} establishes that conservation is the defining
property of the primitive itself --- it is what \emph{is} the count --- and identifies the
Q-event as the indivisible step of the count. Since Actualization is the process that
applies the count (Theorem~\ref{c03:thm:process-face}), the conservation law is inherited:
the process conserves exactly what the primitive defines. No separate conservation postulate
is required.
\end{proof}

\section{The \texorpdfstring{$N=96$}{N=96} Attractor}
\label{c03:sec:n96-attractor}

The \emph{closure} of the theory is the stable fixed point of the actualization process: the
configuration at which the count-producing dynamics has converged, with zero residual
change. The boundary of the theory is this fixed point --- the topology saturates at zero
residual link growth, and every initial pattern converges to the same final geometry. This
is the fixed point that Chapter~\ref{c01:chap:difference} invoked in characterizing (not
deriving) the fundamental boundary (Corollary~\ref{c01:cor:closure-characterizes}); here we
see its origin: the fixed point is the convergence of Difference's own count-producing
process.

\begin{theorem}[The attractor is \texorpdfstring{$N=96$}{N=96}]
\label{c03:thm:n96}
The converged attractor of the actualization process is the $N=96$ network: the size $N=96$
is the stable fixed point of the dynamics, not a freely chosen input.
\end{theorem}

\begin{proof}
The boundary-origin audit (the Closure Principle \cite{c03:QG282}) establishes that the
observable sector is the converged attractor of the actualization dynamics. The dynamics has
a self-reinforcing feedback (activity to links to activity) bounded by network capacity, so
positive feedback saturates at a stable fixed point: zero residual link growth and a
content-independent final geometry. The convergence to the $N=96$ attractor follows from the
closure principle \cite{c03:QG282} and the reduction program \cite{c03:QG293,c03:QG294};
$N=96$ is thereby an output of the dynamics --- the inevitable fixed point --- not a
primitive or a free parameter.
\end{proof}

\begin{corollary}[The spectrum is the attractor's eigenspectrum]
\label{c03:cor:spectrum-eigenspectrum}
The D96 spectrum is the Laplacian eigenspectrum of the converged $N=96$ network: the
inevitable output of the actualization attractor, not an additional primitive.
\end{corollary}

\begin{proof}
The converged network defines a graph Laplacian; its eigenspectrum is the D96 spectrum, with
the 95 positive modes and the single zero mode (the background). Since the network is the
converged fixed point of the dynamics (Theorem~\ref{c03:thm:n96}), the spectrum is a
deterministic function of the attractor --- an output, not an input. This is the inevitable-
spectrum result of the canonical program \cite{c03:QG293,c03:QG294}.
\end{proof}

\section{Actualization and Resonance}
\label{c03:sec:actualization-resonance}

\begin{theorem}[Resonance is actualization's readout]
\label{c03:thm:resonance-readout}
\emph{Resonance} --- the spectral readout of the actualization process, the operator layer
projecting the converged spectrum onto the observable structure --- is actualization's
readout: the sector structure of observable physics is a projection of the one
actualization-driven spectrum. The resonance condition of the theory is the conjunction of
count conservation (the definitional identity) and the closure boundary (the fixed point):
Resonance equals Conservation plus Boundary.
\end{theorem}

\begin{proof}
The sector-emergence audit \cite{c03:QG272} establishes that distinct sectors --- masses,
couplings, mixings, gravity, cosmology --- are projection classes of a single operator layer
over the one spectrum, with no sector-exclusive operator. Since the spectrum is the output
of the actualization attractor (Corollary~\ref{c03:cor:spectrum-eigenspectrum}), the resonance
structure read from it is the readout of the actualization process: one dynamics, one
spectrum, one operator basis, projected onto distinct physical questions. Conservation is
the definitional identity of the count (Theorem~\ref{c03:thm:count-conservation}); the
boundary is the closure fixed point $N=96$ (Theorem~\ref{c03:thm:n96}); the spectrum is
determined by the process that conserves its count and closes at its boundary. The resonance
condition is therefore the conjunction of the two --- Conservation plus Boundary ---
actualization's spectral face.
\end{proof}

\section{Actualization as Derived Necessity}
\label{c03:sec:derived-necessity}

\begin{theorem}[Actualization is derived, not primitive]
\label{c03:thm:derived-not-primitive}
Actualization is INDISPENSABLE as a derivation layer --- without it there is no count, no
network, no spectrum, and no physics --- but it is a DERIVED necessity: indispensable as a
derivation step, yet fully dependent on the primitive Difference and therefore not an
underivable input.
\end{theorem}

\begin{proof}
The hierarchy-necessity audit \cite{c03:QG293} classifies Actualization as INDISPENSABLE: it is
the count-producing process, so without it there is no count, no network, no spectrum, and
no physics. But necessity as a step is not primitiveness as an input: by
Theorem~\ref{c03:thm:process-face}, Actualization is the process face of Difference; the
foundation stress test \cite{c03:QG292} establishes that removing Difference collapses
Actualization (it needs the definitional unit) while removing $\eta$ leaves it intact
(counting needs no metric reference). The minimal-theory audit \cite{c03:QG294} confirms that
the minimal hierarchy Difference → Actualization → Spectrum → Physics is complete, with
Actualization an essential derived step; the MONO006 resolution \cite{c03:MONO006} concludes:
Actualization is derived from Difference alone, indispensable as a step, not a third
primitive.
\end{proof}

The canonical architecture is therefore unchanged --- primitives $\{\text{Difference},\eta\}$,
Actualization derived from Difference --- consistent with Chapters~\ref{c01:chap:difference}
and \ref{c02:chap:eta}; consistent with MONO007 \cite{c03:MONO007}, all statements that could
be read as ascribing primitive status to Actualization are worded to make the derivation
explicit, and no ambiguity remains.

\section{Consequences}
\label{c03:sec:actualization-consequences}

\begin{theorem}[The spectrum is the inevitable output]
\label{c03:thm:spectrum-inevitable}
The D96 spectrum is the inevitable output of the actualization attractor: a deterministic
function of the converged network, carrying no independent choice.
\end{theorem}

\begin{proof}
By Corollary~\ref{c03:cor:spectrum-eigenspectrum}, the spectrum is the Laplacian eigenspectrum of
the converged $N=96$ attractor. The attractor is unique and content-independent
(Theorem~\ref{c03:thm:n96}); the eigenspectrum is a deterministic function of it. Hence the
spectrum carries no choice and is inevitable --- the output of the actualization process,
not an input to it. The reduction program over the fuller hierarchy thereby closes at the
minimal hierarchy Difference → Actualization → Spectrum → Physics: the hierarchy-necessity
audit \cite{c03:QG293} classifies the intermediate layers --- Closure COMPRESSIBLE into
Actualization (the fixed point of the process), Conservation COMPRESSIBLE into
Difference/network (the definitional identity plus the graph handshake), Resonance
COMPRESSIBLE into Spectrum (the operator readout) --- and the minimal-theory audit
\cite{c03:QG294} verifies that every compressed layer is derivable from the minimal hierarchy,
with none actually required as an input.
\end{proof}

\begin{corollary}[Physics is the readout of the actualized spectrum]
\label{c03:cor:physics-readout}
All observable physics is the readout of the actualized spectrum: the resonance structure
projected by the operator layer, from which the sectors of physics emerge.
\end{corollary}

\section{Summary}
\label{c03:sec:actualization-summary}

This chapter has established the second member of the canonical architecture. Actualization
is the process face of Difference (Theorem~\ref{c03:thm:process-face}), whose count
conservation is the primitive's definitional identity (Theorem~\ref{c03:thm:count-conservation})
and whose dynamics converge to the stable fixed point $N=96$
(Theorem~\ref{c03:thm:n96}, Corollary~\ref{c03:cor:spectrum-eigenspectrum}); the resonance
structure read from the resulting spectrum satisfies Resonance equals Conservation plus
Boundary (Theorem~\ref{c03:thm:resonance-readout}). Actualization is indispensable as a layer
yet derived, not primitive (Theorem~\ref{c03:thm:derived-not-primitive}), and the reduction
closes at the minimal hierarchy, with physics the readout of the actualized spectrum
(Theorem~\ref{c03:thm:spectrum-inevitable}, Corollary~\ref{c03:cor:physics-readout}).

The following chapter derives, from this process, the spectrum itself.

\begin{thebibliography}{99}
\bibitem{c03:QG268} AT-QG Phase 268, \emph{Count Conservation Origin}, Docs/Research.
\bibitem{c03:QG272} AT-QG Phase 272, \emph{Sector Emergence Audit}, Docs/Research.
\bibitem{c03:QG282} AT-QG Phase 282, \emph{Boundary Origin Audit} (Closure Principle), Docs/Research.
\bibitem{c03:QG292} AT-QG Phase 292, \emph{Foundation Stress Test}, Docs/Research.
\bibitem{c03:QG293} AT-QG Phase 293, \emph{Hierarchy Necessity Audit}, Docs/Research.
\bibitem{c03:QG294} AT-QG Phase 294, \emph{Minimal Theory Audit}, Docs/Research.
\bibitem{c03:MONO006} AT-MONO006, \emph{A01 Actualization Origin Audit}, Docs/Zenodo/Monograph.
\bibitem{c03:MONO007} AT-MONO007, \emph{Referee-Readiness Resolution}, Docs/Zenodo/Monograph.
\end{thebibliography} after the preamble that defines
%         the theorem environments below (or include via the master file).
% ======================================================================

% ---- Theorem environments (define in the master preamble if not present) ----
% \newtheorem{definition}{Definition}[chapter]
% \newtheorem{lemma}[definition]{Lemma}
% \newtheorem{theorem}[definition]{Theorem}
% \newtheorem{corollary}[definition]{Corollary}
% \newtheorem{remark}[definition]{Remark}

\chapter{Actualization}
\label{c03:chap:actualization}

\section{Introduction}
\label{c03:sec:actualization-introduction}

In Chapter~\ref{c01:chap:difference} we established that Difference is the fundamental boundary
of the theory --- the primitive whose count-conservation law is its definitional identity ---
and in Chapter~\ref{c02:chap:eta} that the tensor reference $\eta$ is the second primitive,
against which the tensor face of Difference is read. The canonical hierarchy
\eqref{c03:eq:core} then proceeds from the primitive pair to the inevitable spectrum through a
process:

\begin{equation}
\label{c03:eq:core}
\text{Difference}\;\longrightarrow\;\text{Actualization}\;\longrightarrow\;\text{Inevitable Spectrum}\;\longrightarrow\;\text{Physics}.
\end{equation}

This chapter develops the second member of that hierarchy: \emph{Actualization}. We
establish, with formal precision, that Actualization is the \emph{process face of
Difference} --- the count-producing dynamics whose definitional identity is count
conservation \cite{c03:QG268}; that its dynamics converge to the stable fixed point $N=96$, the
closure of the theory \cite{c03:QG282}; that the resonance structure of the spectrum is
actualization's readout \cite{c03:QG272}; and --- decisively, consistent with the MONO006
resolution \cite{c03:MONO006} --- that Actualization is a \emph{derived necessity}, not a third
primitive.

\section{Actualization as the Process Face of Difference}
\label{c03:sec:process-face}

\begin{theorem}[Actualization is the process face of Difference]
\label{c03:thm:process-face}
\emph{Actualization} --- the count-producing dynamics by which a Q-event, the unit of
Difference \cite{c03:QG268}, is applied, generating the network whose convergence yields the
spectrum --- is not an independent entity: it is the process face of the one primitive
Difference. Consequently the primitives of the theory are exactly
$\{\text{Difference},\eta\}$: Actualization is their derived process face, not a primitive.
\end{theorem}

\begin{proof}
Count conservation is the definitional identity of the primitive \cite{c03:QG268}: Difference is
the counting difference from a uniform background, and a Q-event is its unit. Actualization
is precisely the process that applies those units, and its identity is inherited from the
primitive: there is no Actualization without the unit, and no unit without Difference. The
minimal primitive set was established as the pair $\{\text{Difference},\eta\}$
(Theorem~\ref{c02:thm:minimal-foundation}, Chapter~\ref{c02:chap:eta}); the MONO006 resolution
\cite{c03:MONO006} confirms this classification, and the MONO007 architecture \cite{c03:MONO007}
presents Actualization as a derived layer. Hence there is no primitive-status ambiguity:
Actualization is the dynamical face of Difference, not a separate primitive.
\end{proof}

\section{Count Conservation}
\label{c03:sec:count-conservation}

\begin{theorem}[Count conservation is the definitional identity]
\label{c03:thm:count-conservation}
\emph{Count conservation} --- the total count of units of Difference is conserved by the
actualization process: Q-events are neither created nor destroyed, only redistributed ---
is not a separate postulate but the definitional identity of the primitive Difference and
therefore of its process face Actualization. One application of Actualization produces one
unit of count, so conservation follows automatically from the primitive, not as an imported
assumption of the theory.
\end{theorem}

\begin{proof}
The count-conservation origin \cite{c03:QG268} establishes that conservation is the defining
property of the primitive itself --- it is what \emph{is} the count --- and identifies the
Q-event as the indivisible step of the count. Since Actualization is the process that
applies the count (Theorem~\ref{c03:thm:process-face}), the conservation law is inherited:
the process conserves exactly what the primitive defines. No separate conservation postulate
is required.
\end{proof}

\section{The \texorpdfstring{$N=96$}{N=96} Attractor}
\label{c03:sec:n96-attractor}

The \emph{closure} of the theory is the stable fixed point of the actualization process: the
configuration at which the count-producing dynamics has converged, with zero residual
change. The boundary of the theory is this fixed point --- the topology saturates at zero
residual link growth, and every initial pattern converges to the same final geometry. This
is the fixed point that Chapter~\ref{c01:chap:difference} invoked in characterizing (not
deriving) the fundamental boundary (Corollary~\ref{c01:cor:closure-characterizes}); here we
see its origin: the fixed point is the convergence of Difference's own count-producing
process.

\begin{theorem}[The attractor is \texorpdfstring{$N=96$}{N=96}]
\label{c03:thm:n96}
The converged attractor of the actualization process is the $N=96$ network: the size $N=96$
is the stable fixed point of the dynamics, not a freely chosen input.
\end{theorem}

\begin{proof}
The boundary-origin audit (the Closure Principle \cite{c03:QG282}) establishes that the
observable sector is the converged attractor of the actualization dynamics. The dynamics has
a self-reinforcing feedback (activity to links to activity) bounded by network capacity, so
positive feedback saturates at a stable fixed point: zero residual link growth and a
content-independent final geometry. The convergence to the $N=96$ attractor follows from the
closure principle \cite{c03:QG282} and the reduction program \cite{c03:QG293,c03:QG294};
$N=96$ is thereby an output of the dynamics --- the inevitable fixed point --- not a
primitive or a free parameter.
\end{proof}

\begin{corollary}[The spectrum is the attractor's eigenspectrum]
\label{c03:cor:spectrum-eigenspectrum}
The D96 spectrum is the Laplacian eigenspectrum of the converged $N=96$ network: the
inevitable output of the actualization attractor, not an additional primitive.
\end{corollary}

\begin{proof}
The converged network defines a graph Laplacian; its eigenspectrum is the D96 spectrum, with
the 95 positive modes and the single zero mode (the background). Since the network is the
converged fixed point of the dynamics (Theorem~\ref{c03:thm:n96}), the spectrum is a
deterministic function of the attractor --- an output, not an input. This is the inevitable-
spectrum result of the canonical program \cite{c03:QG293,c03:QG294}.
\end{proof}

\section{Actualization and Resonance}
\label{c03:sec:actualization-resonance}

\begin{theorem}[Resonance is actualization's readout]
\label{c03:thm:resonance-readout}
\emph{Resonance} --- the spectral readout of the actualization process, the operator layer
projecting the converged spectrum onto the observable structure --- is actualization's
readout: the sector structure of observable physics is a projection of the one
actualization-driven spectrum. The resonance condition of the theory is the conjunction of
count conservation (the definitional identity) and the closure boundary (the fixed point):
Resonance equals Conservation plus Boundary.
\end{theorem}

\begin{proof}
The sector-emergence audit \cite{c03:QG272} establishes that distinct sectors --- masses,
couplings, mixings, gravity, cosmology --- are projection classes of a single operator layer
over the one spectrum, with no sector-exclusive operator. Since the spectrum is the output
of the actualization attractor (Corollary~\ref{c03:cor:spectrum-eigenspectrum}), the resonance
structure read from it is the readout of the actualization process: one dynamics, one
spectrum, one operator basis, projected onto distinct physical questions. Conservation is
the definitional identity of the count (Theorem~\ref{c03:thm:count-conservation}); the
boundary is the closure fixed point $N=96$ (Theorem~\ref{c03:thm:n96}); the spectrum is
determined by the process that conserves its count and closes at its boundary. The resonance
condition is therefore the conjunction of the two --- Conservation plus Boundary ---
actualization's spectral face.
\end{proof}

\section{Actualization as Derived Necessity}
\label{c03:sec:derived-necessity}

\begin{theorem}[Actualization is derived, not primitive]
\label{c03:thm:derived-not-primitive}
Actualization is INDISPENSABLE as a derivation layer --- without it there is no count, no
network, no spectrum, and no physics --- but it is a DERIVED necessity: indispensable as a
derivation step, yet fully dependent on the primitive Difference and therefore not an
underivable input.
\end{theorem}

\begin{proof}
The hierarchy-necessity audit \cite{c03:QG293} classifies Actualization as INDISPENSABLE: it is
the count-producing process, so without it there is no count, no network, no spectrum, and
no physics. But necessity as a step is not primitiveness as an input: by
Theorem~\ref{c03:thm:process-face}, Actualization is the process face of Difference; the
foundation stress test \cite{c03:QG292} establishes that removing Difference collapses
Actualization (it needs the definitional unit) while removing $\eta$ leaves it intact
(counting needs no metric reference). The minimal-theory audit \cite{c03:QG294} confirms that
the minimal hierarchy Difference → Actualization → Spectrum → Physics is complete, with
Actualization an essential derived step; the MONO006 resolution \cite{c03:MONO006} concludes:
Actualization is derived from Difference alone, indispensable as a step, not a third
primitive.
\end{proof}

The canonical architecture is therefore unchanged --- primitives $\{\text{Difference},\eta\}$,
Actualization derived from Difference --- consistent with Chapters~\ref{c01:chap:difference}
and \ref{c02:chap:eta}; consistent with MONO007 \cite{c03:MONO007}, all statements that could
be read as ascribing primitive status to Actualization are worded to make the derivation
explicit, and no ambiguity remains.

\section{Consequences}
\label{c03:sec:actualization-consequences}

\begin{theorem}[The spectrum is the inevitable output]
\label{c03:thm:spectrum-inevitable}
The D96 spectrum is the inevitable output of the actualization attractor: a deterministic
function of the converged network, carrying no independent choice.
\end{theorem}

\begin{proof}
By Corollary~\ref{c03:cor:spectrum-eigenspectrum}, the spectrum is the Laplacian eigenspectrum of
the converged $N=96$ attractor. The attractor is unique and content-independent
(Theorem~\ref{c03:thm:n96}); the eigenspectrum is a deterministic function of it. Hence the
spectrum carries no choice and is inevitable --- the output of the actualization process,
not an input to it. The reduction program over the fuller hierarchy thereby closes at the
minimal hierarchy Difference → Actualization → Spectrum → Physics: the hierarchy-necessity
audit \cite{c03:QG293} classifies the intermediate layers --- Closure COMPRESSIBLE into
Actualization (the fixed point of the process), Conservation COMPRESSIBLE into
Difference/network (the definitional identity plus the graph handshake), Resonance
COMPRESSIBLE into Spectrum (the operator readout) --- and the minimal-theory audit
\cite{c03:QG294} verifies that every compressed layer is derivable from the minimal hierarchy,
with none actually required as an input.
\end{proof}

\begin{corollary}[Physics is the readout of the actualized spectrum]
\label{c03:cor:physics-readout}
All observable physics is the readout of the actualized spectrum: the resonance structure
projected by the operator layer, from which the sectors of physics emerge.
\end{corollary}

\section{Summary}
\label{c03:sec:actualization-summary}

This chapter has established the second member of the canonical architecture. Actualization
is the process face of Difference (Theorem~\ref{c03:thm:process-face}), whose count
conservation is the primitive's definitional identity (Theorem~\ref{c03:thm:count-conservation})
and whose dynamics converge to the stable fixed point $N=96$
(Theorem~\ref{c03:thm:n96}, Corollary~\ref{c03:cor:spectrum-eigenspectrum}); the resonance
structure read from the resulting spectrum satisfies Resonance equals Conservation plus
Boundary (Theorem~\ref{c03:thm:resonance-readout}). Actualization is indispensable as a layer
yet derived, not primitive (Theorem~\ref{c03:thm:derived-not-primitive}), and the reduction
closes at the minimal hierarchy, with physics the readout of the actualized spectrum
(Theorem~\ref{c03:thm:spectrum-inevitable}, Corollary~\ref{c03:cor:physics-readout}).

The following chapter derives, from this process, the spectrum itself.

\begin{thebibliography}{99}
\bibitem{c03:QG268} AT-QG Phase 268, \emph{Count Conservation Origin}, Docs/Research.
\bibitem{c03:QG272} AT-QG Phase 272, \emph{Sector Emergence Audit}, Docs/Research.
\bibitem{c03:QG282} AT-QG Phase 282, \emph{Boundary Origin Audit} (Closure Principle), Docs/Research.
\bibitem{c03:QG292} AT-QG Phase 292, \emph{Foundation Stress Test}, Docs/Research.
\bibitem{c03:QG293} AT-QG Phase 293, \emph{Hierarchy Necessity Audit}, Docs/Research.
\bibitem{c03:QG294} AT-QG Phase 294, \emph{Minimal Theory Audit}, Docs/Research.
\bibitem{c03:MONO006} AT-MONO006, \emph{A01 Actualization Origin Audit}, Docs/Zenodo/Monograph.
\bibitem{c03:MONO007} AT-MONO007, \emph{Referee-Readiness Resolution}, Docs/Zenodo/Monograph.
\end{thebibliography}